\section{Introduction}
\label{sec:introduction}

A sequential experiment may need to retain several promising alternatives
while new alternatives continue to enter. Examples include a testing platform
that receives externally generated designs and an experimental pipeline that
introduces candidates during an ongoing study. The investigator must decide
which alternative to sample and what the accumulated evidence says about the
whole candidate set. Cumulative reward loss measures the quality of allocation,
but a small loss does not guarantee that all alternatives within a specified
tolerance of the current best can be distinguished from those outside it.

Consider a collection of well-studied candidates and one newly available arm.
Although the earlier means have not changed, uncertainty about the new arm's
possible superiority may prevent positive certification of many earlier
candidates. An informative upper bound on its mean may be needed before the
earlier candidates can be certified. A candidate close to today's relative
threshold may itself require many observations to classify. After a sufficiently
better arm arrives, both possible sides of
the old threshold can lie well below the new one. Observations already collected
can then suffice to exclude the candidate, without resolving its earlier
membership. Continual reporting therefore creates shared waiting costs, while
later arrivals can make unfinished fine-resolution tasks unnecessary.

We study these effects in the stationary stochastic bandit model with exogenous
permanent arrivals. Regret is measured against the best arm available at each
physical observation time. This model and comparator are established in
ballooning bandits \citep{ghalme2021}. For any finite prefix, permanent arrival
is also a monotone availability pattern in the sleeping-bandit formulation
\citep{kleinberg2010}. An ordinary index policy can admit an arm on arrival.
The additional statistical requirement considered here is a reliable report of
the complete current relative near-optimal set, together with a bound on how
long its members and nonmembers remain unresolved.

The fixed-arm version of relative-good-set identification is the
ALL-$\varepsilon$ problem of \citet{mason2020}. Their analysis already
distinguishes discrimination near the membership boundary from uncertainty
about the best mean. Their algorithms, including $(\mathrm{ST})^2$ and FAREAST,
are essential comparisons for any proposed extension. Reusing confidence bounds
after an arrival is a natural adaptation; the cost of maintaining their reports
through repeated changes of the relative threshold requires further analysis.
Our lower bounds measure cumulative unresolved arm time under one simultaneous
correctness requirement. They concern continuing reports, rather than a single
terminal recommendation or independently imposed deadlines.
Fixed-budget thresholding also studies terminal classification loss
\citep{cheshire2020}; cumulative unresolved time records the full reporting path.

Joint reward and identification objectives also have substantial precedents.
\citet{degenne2019} study the compromise between regret and decision time;
\citet{yang2026} characterize regret-optimal fixed-confidence best-arm
identification and its sampling cost. \citet{cho2025} combine reward-oriented
sampling with anytime-valid testing for fixed-threshold good-arm
identification. Separating a reward policy from a reporting layer thus has
established precedents. Sequential multiple testing
and support recovery likewise supply the simultaneous-recovery scales used
below \citep{malloy2014,jamieson2018}. We study how a continuing obligation to
report the relative set imposes these costs, shares them among arms, and
sometimes allows them to be avoided.

Our first main result gives a common-environment lower bound: an insufficiently
studied arrival can keep a whole cohort unresolved. Combined with an executable
upper bound, it identifies the polynomial arrival exponent below which both
expected regret and cumulative unresolved area can be sublinear. A separate
worst-schedule result characterizes the area as the smaller of the available
arm-time and a quadratic-in-arm-count identification cost. Its lower bound
already occurs for an initial batch. Repeated benchmark improvements have a
further temporal effect, which we analyze separately.

Our second main result isolates the cost of temporary membership questions.
In an explicit multiple-update family, a small area budget forces enough
fine-resolution classifications before they expire. At a specified sequence
of horizons this gives matching rapid-maintenance regret bounds. A larger
constant-factor area budget permits coarse evidence to be reused after the
next improvement and has a substantially smaller attainable regret. We also
prove a fixed-probability regret lower bound under precisely the same area
requirement. These results describe two service budgets; they do not identify
the entire instancewise frontier.

The checkpoint construction is deliberately simple. It accepts arms immediately,
retains all observations, and interleaves reporting service with full-history
UCB. It provides a constructive bound whose reward guarantee holds without
conditioning on correct reports. Its limitations are visible in the numerical
comparisons: strong confidence-based adaptations can improve both regret and
area, and fine fixed checkpoints can be impractical at moderate budgets.
The construction establishes attainable joint costs and explains their
mechanism; it need not be the best practical implementation.
